\section*{Wstęp}
\addcontentsline{toc}{section}{Wstęp}
\label{chap:Wstęp}
\markboth{WSTĘP}{WSTĘP}

Powodzie stanowią jedno z~najpoważniejszych zagrożeń naturalnych dla społeczeństw i~infrastruktury na całym świecie. Według raportu Światowej Organizacji Meteorologicznej~\cite{wmo2021}, w~ciągu ostatnich dwóch dekad liczba katastrof związanych z~wodą wzrosła o~70\%, przy czym powodzie odpowiadają za największą liczbę ofiar śmiertelnych i~strat ekonomicznych wśród wszystkich zjawisk hydrometeorologicznych. W~kontekście zmian klimatycznych problem ten narasta -- zwiększająca się częstotliwość ekstremalnych opadów deszczu oraz rosnący poziom mórz prowadzą do wzrostu ryzyka powodziowego w~wielu regionach świata, w~tym w~Polsce i~Europie~\cite{undrr2022}.

Szczególnie podatna na zagrożenie powodziowe jest infrastruktura krytyczna, obejmująca szpitale, szkoły, drogi, mosty oraz inne obiekty użyteczności publicznej~\cite{koks2019}. Zatopienie czy uszkodzenie takich obiektów może prowadzić do kaskadowych skutków dla całych społeczności -- od braku dostępu do opieki medycznej, przez izolację komunikacyjną, po długotrwałe zakłócenia w~funkcjonowaniu administracji lokalnej. Identyfikacja i~kwantyfikacja ryzyka powodziowego dla infrastruktury krytycznej stanowi zatem kluczowy element zarządzania ryzykiem katastroficznym oraz planowania przestrzennego~\cite{chukwuma2021flood}.

Tradycyjne podejścia do oceny ryzyka powodziowego opierają się na specjalistycznych systemach Geograficznych Systemów Informacyjnych (GIS), które wymagają wysokich kompetencji technicznych oraz znajomości skomplikowanych narzędzi analitycznych. Stanowi to istotną barierę dla potencjalnych użytkowników, szczególnie w~kontekście Rapid Mapping -- szybkiego mapowania sytuacji kryzysowej wymagającego wydania oceny ryzyka w~ciągu godzin od zdarzenia. Służby reagowania kryzysowego, decydenci lokalni czy planiści przestrzenni często nie dysponują czasem ani zasobami na przeprowadzenie złożonych analiz geoprzestrzennych przy użyciu tradycyjnych narzędzi GIS.

W~ostatnich latach rozwinęły się dwa kierunki technologiczne istotne dla tego problemu. Po pierwsze, metody uczenia maszynowego umożliwiają automatyzację przetwarzania danych satelitarnych i~budowę modeli predykcyjnych ryzyka powodziowego~\cite{waleed2025advancing}. Po drugie, duże modele językowe (Large Language Models, LLM) i~agenci AI pozwalają na tworzenie interfejsów języka naturalnego, dzięki którym użytkownicy mogą zadawać pytania bez znajomości składni GIS czy języków programowania~\cite{li2024chatgeoai}.

Jednocześnie rozwijają się otwarte źródła danych geoprzestrzennych. Program Copernicus Unii Europejskiej oferuje globalny system monitorowania powodzi (GloFAS) dostarczający prognoz i~danych o~zagrożeniu powodziowym w~czasie rzeczywistym~\cite{jrc2020glofas}. Overture Maps Foundation, inicjatywa wspierana przez wiodące firmy technologiczne (Amazon, Meta, Microsoft, TomTom), udostępnia dane o~infrastrukturze globalnej, obejmujące 2,6 miliarda budynków, 64 miliony punktów użyteczności publicznej oraz 321 milionów segmentów drogowych~\cite{overture2024foundation}.

Praca łączy te technologie w~prototypowym systemie RiskScoop AI -- agencie AI umożliwiającym analizę ryzyka powodziowego dla infrastruktury krytycznej poprzez zapytania w~języku naturalnym. Na rynku brakuje narzędzi łączących interfejs języka naturalnego z~oceną ryzyka powodziowego dla infrastruktury krytycznej, szczególnie w~kontekście Rapid Mapping. Rozwiązanie bazuje na osiągnięciach z~dziedziny agentów geoprzestrzennych (ChatGeoAI, LLM-Find, GeoGPT) oraz otwartych europejskich źródłach danych (Copernicus GloFAS, Overture Maps).

Praca ma następującą strukturę. Rozdział pierwszy przedstawia przegląd literatury obejmujący agentów AI w~analizie geoprzestrzennej, systemy monitorowania powodzi i~otwarte dane geoprzestrzenne, a~następnie analizę istniejących rozwiązań komercyjnych i~akademickich wraz z~identyfikacją luki rynkowej oraz wymagania funkcjonalne i~niefunkcjonalne systemu. Rozdział drugi prezentuje projekt systemu, opisując architekturę, narzędzia agenta oraz model ryzyka. Rozdział trzeci omawia szczegóły implementacji backendu i~frontendu. Rozdział czwarty przedstawia przeprowadzone testy funkcjonalne, wydajnościowe oraz testy interpretacji języka naturalnego wraz z~analizą wyników. Pracę zamyka podsumowanie oraz wnioski dotyczące osiągniętych rezultatów i~kierunków dalszego rozwoju systemu.

Prototyp stanowi dowód koncepcji (proof of concept) -- użytkownik może zadać pytanie typu ,,Które szpitale w~Warszawie są zagrożone powodzią?'' i~otrzymać wynik na mapie bez znajomości narzędzi GIS. System działa w~przeglądarce i~korzysta wyłącznie z~otwartych źródeł danych.
